\documentclass[12pt, letterpaper]{article}
\usepackage{graphicx} % Required for inserting images
%\usepackage[T1]{fontenc}

\title{Programy użytkowe \\ Lista 10}
\date{}
\usepackage{polski}
\begin{document}
\maketitle
\begin{enumerate}
\item Utwórz najprostszy plik z kodem w \LaTeX{u} na dysku komputera. Używając w terminalu polecenia \verb|latexmk| skompiluj ten plik do pliku pdf. Uwaga: podczas kompilacji powstaje wiele plików pośrednich. Zapoznaj się z opcjami \verb|-output-directory| oraz \verb|-c| (aby umieścić pliki w osobnym katalogu oraz posprzątać po kompilacji).
\item Napisz skrypt, który będzie wykonywał polecenia z poprzedniego zadania.
\item Napisz skrypt, który zapisuje w pliku kod \LaTeX{a}, służący do wyświetlenia podanego wykresu. Wskazówka: użyj operatora zapisu do pliku \verb|>>|.
\item Napisz jeden skrypt, który umieści w pliku \LaTeX{a} wykresy generowane przez skrypt z poprzedniej listy wraz z odpowiednimi podpisami, skompiluje ten plik oraz usunie pliki dodatkowe. Skrypt ma przyjmować liczbę wykresów do wygenerowania jako argument wiersza poleceń.
\item Zaponaj się z działaniem programów \verb|latexindent| oraz \verb|latexdiff| i je samodzielnie wypróbuj.
\end{enumerate}
\end{document}